% =====================================================================
%  TRABAJOS FUTUROS (versión recortada, para comparar con
%  capitulos/10-trabajos-futuros.tex)
%
%  Revisión del 27 de octubre: las cinco líneas que pidieron los autores
%  (navegación, cómputo, evasión de obstáculos, otro controlador e interacción
%  por voz). Las seis anteriores, de campañas más grandes a evaluación con
%  usuarios, se quitaron por indicación de los autores.
% =====================================================================
\chapter*{Trabajos futuros}
\addcontentsline{toc}{chapter}{TRABAJOS FUTUROS}

\paragraph{Una navegación más precisa.} La navegación de los agentes se puede
mejorar en dos frentes. El primero es la tracción: el agente no se mueve por
debajo de 0,40~m/s, y calibrar el variador o cambiar la relación de transmisión
le daría el régimen lento que hoy le falta para la aproximación final. El
segundo es la medida del avance: codificadores en las ruedas, u odometría visual
con las cámaras que el agente ya lleva, darían una medida que no depende de que
el pasillo tenga estructura. Con las dos, la tolerancia de llegada podría volver
a 0,25~m y el sistema dejaría de depender del entorno para localizarse.

\paragraph{Más capacidad de cómputo.} La tarjeta original, de dos núcleos,
queda entre el 54 y el 59~\% de uso en reposo con la navegación completa, y su
carga se dispara cuando los procesos de los dos agentes se anuncian en la red.
Adaptar la plataforma a una Raspberry Pi~4 o~5, de cuatro núcleos, daría margen
para la navegación, para un controlador más exigente y para la evasión de
obstáculos. Para la Raspberry Pi~4, la comunidad del DeepRacer ya publica el
software del vehículo.

\paragraph{Evasión de obstáculos.} Las pruebas se hicieron con los pasillos
despejados. En un edificio en uso, el agente tiene que esquivar a las personas
que se cruzan y los objetos que cambian de sitio, y no solo detenerse ante
ellos. Integrar la evasión de obstáculos en movimiento, y probarla con
personas que caminan por el pasillo, es el paso previo a operar en horario de
clases.

\paragraph{Un controlador predictivo.} El seguimiento de trayectorias usa
Regulated Pure Pursuit, que sigue la ruta punto a punto y cuya guarda de
seguridad anula la curvatura cuando una inversión de marcha queda cerca de la
meta. Un controlador predictivo, como el MPPI que incluye Nav2 y que admite
cinemática tipo automóvil, optimiza el movimiento varios pasos hacia adelante y
tiene en cuenta los obstáculos. Con él, la media vuelta podría volver a Nav2 y
la aproximación final sería más precisa. Exige más cómputo, y por eso va de la
mano de la línea anterior.

\paragraph{Interacción por voz.} La interfaz pide elegir el origen y el destino
en listas. Con reconocimiento de voz a texto y un modelo de lenguaje, el usuario
podría pedir el servicio hablando, por ejemplo «lléveme al Salón 402», y el
sistema le explicaría cada etapa en voz alta. Como la red del edificio puede no
tener salida a internet, los dos modelos tendrían que correr en el computador del coordinador. El de lenguaje solo traduciría la petición a un origen y un
destino del catálogo, sin decidir el recorrido, que sigue a cargo del
planificador.
